{\color{red}WIP}\\

In this section, I will outline the relationship between the following components and the practical steps needed to establish the relationship:

\begin{enumerate}
    \item Artha - Meaningful Programming Language
    \item DAL - Design Abstraction Language
    \item CSM - Computable Semantic Model
    \item CSP - Computable Semantic Primitives
    \item Engine - Semantic Reasoning Engine
    \item CLP - Compressed Log Processor
    \item SI - Synthesized Implementation
\end{enumerate}

When discussing this framework, the core principles revolve around meaning. So I think I will start by establishing that meaning is constructed by narratives that use the primitive meaningful building blocks. The narratives themselves have meaning and structure and that structure is used to automatically build a semantic model written in the design abstraction language. Since the meaning was built from computable semantic primitives, which are themselves established using the DAL, the semantic model becomes executable and becomes a computable semantic model.

Once meaning has been established, now it can be used in another narrative. It is a compositional system where meaning is composed through narratives and used to compose more meaning. Since the meaning is built from computable semantic primitives and it has an unambiguous structure, it can be synthesized into an implementation. In the process, the meaning of the software system has been automatically developed in a programming language.

The engine reasons about the meaning using the structure established by the computable semantic model. Since the structure is unambiguous and meaningful, it is able to reason about it to enforce its own definition of correctness and to also identify missing or incorrect meaning. Through the synthesis and instrumentation, the necessary information needed to replay the execution is preserved using domain specific compression and the engine is able to diagnose any failures and enable the meaning to expand through root cause analysis. This process automatically manages the software system and enables a deterministic learning loop. 

Since the meaning of the participant is specified as part of the narrative, its domain structure can also be included in that meaning. So the meaning is what CLP uses to perform the domain specific compression, the domain structure is part of its meaning. If you imagine a story, the description of the character (which is part of the meaning) can be as detailed and unambiguous as needed, in this case, that description can be extended to include the domain structure of the participant. 

The meaning can also tell you more about the data (ex. how it changes) so that it can be exploited to optimally compress this. There is a difference between knowing a value is a float and knowing that it is a GPS coordinates latitude or ambient temperature. It also makes the data itself meaningful, for example, if its GPS coordinates, you know when there is movement and when there is no movement. Also, since the meaning is already known, the queries will be meaningful and the ingestion performance can be improved. There are lots of implications for this but I'll discuss them later in the domain specific compression section. Ultimately, I just see meaning as the elimination of ambiguity in the domain (which CLP benefits from) and this framework provides the means to fully eliminate ambiguity in the meaning.

Next, I will discuss how each component can be practically implemented.


